\subsection{Evidencia descriptiva sobre el cumplimiento}

Para el análisis descriptivo se utilizan los resultados de las acciones de fiscalización realizadas por la Sunass durante 2024, el año más reciente con resultados consolidados.\footnote{Las fiscalizaciones efectuadas en 2024 no necesariamente derivaron en sanciones dentro del mismo periodo: algunos procedimientos administrativos sancionadores (PAS) culminaron con resoluciones emitidas en 2024 y otros durante 2025; las fiscalizaciones de 2025 aún no cuentan con la totalidad de los PAS concluidos. Cuando se verificó un incumplimiento, se consideran los resultados definitivos contenidos en las resoluciones sancionadoras, dado que los valores preliminares pueden modificarse como consecuencia de los descargos, medios probatorios o información complementaria presentada por la empresa prestadora. Asimismo, el valor obtenido de una meta puede considerarse igual a cero cuando la información reportada no puede ser validada por la Sunass, como ocurre, por ejemplo, en metas de continuidad o presión que requieren registros técnicos verificables mediante \textit{data logger}.} La Tabla~\ref{tab:cumplimiento_metas_gestion} presenta el cumplimiento promedio observado según el tipo de meta fiscalizada.

\begin{table}[!t]
\centering
\footnotesize
\caption{Cumplimiento promedio de las metas de gestión}
\label{tab:cumplimiento_metas_gestion}
\renewcommand{\arraystretch}{1.15}
\setlength{\tabcolsep}{4pt}

\begin{tabularx}{\columnwidth}{
@{}
>{\raggedright\arraybackslash}X
>{\centering\arraybackslash}p{1.60cm}
@{}
}
\hline
\textbf{Meta de gestión} &
\textbf{Cumplimiento promedio} \\
\hline
Agua no facturada & 97\% \\
Relación de trabajo & 93\% \\
Micromedición & 89\% \\
Continuidad & 83\% \\
Presión & 83\% \\
Catastro comercial & 82\% \\
Catastro técnico & 76\% \\
\hline
Reemplazo de medidores & 68\% \\
Incremento de conexiones activas de agua o alcantarillado & 64\% \\
Instalación de medidores & 64\% \\
Renovación anual de medidores & 55\% \\
\hline
Porcentaje de ejecución de la reserva MRSE & 45\% \\
Porcentaje de ejecución de la reserva GRD y ACC & 42\% \\
Porcentaje de avance financiero del programa de inversiones & 36\% \\
Ejecución de la reserva PCC y otros (PAS, VMA) & 25\% \\
Porcentaje de avance financiero del programa de inversiones del servicio de monitoreo y gestión del uso de aguas subterráneas & 11\% \\
\hline
Sin clasificación & 69\% \\
\hline
\textbf{Total general} & \textbf{71\%} \\
\hline
\end{tabularx}

\vspace{0.12cm}
\begin{minipage}{0.97\columnwidth}
\footnotesize
\textit{Nota:} El cumplimiento promedio corresponde a los resultados observados en las fiscalizaciones realizadas por la Sunass en 2024. La fila ``Sin clasificación'' comprende registros en los que no fue posible identificar el tipo de meta y no se incorpora en la comparación entre categorías.
\end{minipage}
\end{table}

Ordenados de mayor a menor cumplimiento, los indicadores se agrupan descriptivamente según la naturaleza predominante de la variable medida. Esta agrupación no corresponde a la clasificación normativa del SIIGEPSS. Las primeras siete filas de la Tabla~\ref{tab:cumplimiento_metas_gestion} comprenden resultados de la prestación, resultados operativos, una condición económico-operativa y capacidades de gestión; las cuatro siguientes corresponden a ejecución física; y las cinco posteriores, a ejecución financiera. Los registros incluidos en la fila ``Sin clasificación'' se mantienen para calcular el promedio general, pero no se consideran en la comparación entre estos grupos.

El promedio general de cumplimiento asciende a 71\%. Las metas referidas a resultados de la prestación, resultados operativos y determinadas capacidades de gestión se ubican principalmente en la parte superior de la distribución. Las metas de ejecución física y, con mayor claridad, las de ejecución financiera se concentran en la mitad inferior. La diferencia no responde únicamente a un indicador extremo, sino que se observa en distintas metas asociadas con inversiones, reservas y actividades programadas.

Al comparar estos resultados con la frecuencia de asignación presentada en la Tabla~\ref{tab:metas_gestion_estudios_recientes}, se advierte que algunas de las metas con niveles de cumplimiento muy distintos tienen una presencia extendida en los estudios tarifarios. La relación de trabajo y la continuidad, presentes en los 48 estudios analizados, alcanzaron cumplimientos promedio de 93\% y 83\%, respectivamente; el catastro comercial, incluido en 45 estudios, registró un cumplimiento de 82\%. El avance financiero del programa de inversiones, incorporado en 41 de los 48 estudios ---equivalente al 85,4\,\%---, alcanzó en cambio un cumplimiento promedio de 36\%. Las dos tablas corresponden a periodos y universos de observación distintos, por lo que esta comparación no permite estimar una relación estadística entre frecuencia y cumplimiento. Sí permite advertir que el bajo cumplimiento observado en el avance financiero corresponde a una meta de presencia extendida dentro del régimen tarifario y no a un indicador de aplicación marginal.

Dentro de las metas fiscalizadas, nueve de los 16 tipos identificados corresponden a ejecución física o financiera: cuatro a ejecución física y cinco a ejecución financiera. Los siete restantes se distribuyen entre resultados de la prestación, resultados operativos, condiciones económico-operativas y capacidades de gestión. Este conteo no representa el peso efectivo de cada categoría en el ICG de una empresa determinada, porque la combinación de metas varía entre empresas y periodos regulatorios. Sin embargo, muestra que, cuando el índice se calcula como un promedio simple de los ICI, una categoría representada mediante un mayor número de metas puede adquirir una incidencia agregada superior aun cuando la regulación no le asigne explícitamente una ponderación mayor.

Como referencia comparativa, la experiencia internacional examinada en el Anexo~\ref{anexo:experiencia_internacional} muestra otras formas de articular la evaluación del desempeño con el seguimiento de las inversiones. Los marcos revisados difieren en sus mecanismos de evaluación y en las consecuencias asociadas con el desempeño, pero mantienen alguna separación entre la medición de los resultados o condiciones observadas y el control de las inversiones, obras o actividades destinadas a producirlos. La Tabla~\ref{tab:comparacion_internacional_separacion} sintetiza esta distinción.


En el caso peruano, las diferencias de cumplimiento entre tipos de meta no deben interpretarse automáticamente como diferencias de esfuerzo, responsabilidad o desempeño gerencial. Los indicadores operativos pueden depender de decisiones recurrentes, del capital instalado y de las condiciones estructurales de la prestación, mientras que la ejecución de inversiones y reservas también puede estar condicionada por las etapas del ciclo de inversión, los procesos de contratación, la disponibilidad de recursos, las autorizaciones administrativas y la intervención de terceros. La evidencia presentada es descriptiva: no identifica causalmente las razones del menor cumplimiento ni permite estimar la respuesta de las empresas ante una modificación del diseño regulatorio. Sí muestra que un mismo porcentaje de incumplimiento puede transmitir información distinta según la naturaleza del objeto medido y los factores que condicionan su ejecución. Cuando una meta depende de procesos administrativos, financieros o contractuales distintos de aquellos que determinan los resultados de la prestación, su incumplimiento no necesariamente constituye una señal equivalente de desempeño gerencial.

En los resultados agregados del ICG también se observan niveles de cumplimiento reducidos. En las 17 fiscalizaciones del ICG de metas de gestión realizadas durante 2024, el índice alcanzó un promedio de 65,1\% y presentó valores comprendidos entre 10\% y 100\%. Trece fiscalizaciones, equivalentes al 76,5\%, registraron un ICG inferior al umbral de 85\% previsto en la tipificación A--4.1, mientras que solo una alcanzó el cumplimiento integral de las metas incorporadas en el índice.

Dos fiscalizaciones adicionales correspondieron al ICG del servicio de monitoreo y gestión del uso de aguas subterráneas. Sus valores fueron de 10\% y 50\%, con un promedio de 30\%, por lo que ambas se ubicaron por debajo del mismo umbral. Este índice agrega el conjunto de metas correspondientes a dicho servicio.

Antes de generalizar estos resultados, conviene precisar el alcance de las fiscalizaciones observadas. Las metas correspondientes a un año regulatorio se fiscalizan en un momento posterior del ciclo administrativo y no todas las empresas prestadoras tienen asignado el mismo conjunto de metas, incluida la de avance financiero del programa de inversiones. El grupo fiscalizado durante 2024 no constituye, por ello, un universo homogéneo ni equivale al total de empresas reguladas, de modo que estos resultados no permiten estimar la proporción de empresas prestadoras que se encontraría por debajo del 85\,\% en el ámbito nacional. Dentro del grupo fiscalizado, sin embargo, los valores inferiores al umbral fueron frecuentes, y el resultado agregado del ICG puede variar sustancialmente según la composición y el nivel de cumplimiento de las metas evaluadas en cada caso.

Estas cautelas no eliminan la importancia de la composición del ICG para interpretar el desempeño regulatorio. Como el índice se obtiene mediante el promedio de los ICI, la incorporación de metas que presentan distintos niveles de cumplimiento, responden a determinantes no equivalentes y tienen diferente frecuencia de asignación puede modificar la evaluación agregada. Tratar dentro de un mismo promedio los resultados de la prestación, las capacidades de gestión, la ejecución física y la ejecución financiera supone atribuirles una comparabilidad que la evidencia presentada, por sí sola, no permite establecer. Las secciones siguientes examinan, a partir del marco conceptual propuesto, si estos objetos deben incidir de la misma manera en el índice o permanecer sujetos a instrumentos regulatorios diferenciados.